\chapter{Deflation and Nonlinear Eigenproblems}
\label{ch:deflation}

\chapterorient{The last two chapters found a \emph{few} eigenpairs of a giant
pencil. But the engineer rarely wants only the fundamental mode---they want the
fundamental \emph{and} its first several overtones, ten or twenty resonances
peeled off the bottom of the spectrum. How do you stop the iteration from
finding the same lowest mode over and over? The answer is \emph{deflation}: once
a pair has converged, project it out of the search so the iteration is forced to
find the \emph{next} one. That single idea---compute eigenpairs one at a time,
locking each as you go---carries us into the chapter's harder half. When the
cavity is filled with a lossy, dispersive material, the operator itself depends
on the frequency you are solving for: a \emph{nonlinear eigenproblem}
$\mat{T}(\lambda)\vx=\mathbf 0$ in which the resonance feeds back into the
equation that defines it. Linearization can no longer save us. We meet the
deflated quasi-Newton method Palace actually runs---Newton's method on the
eigenvalue, with each found mode deflated through an augmented system---and watch
its inner linear solve split into a tiny dense block and a huge sparse one.}

\section{Why one eigenpair is never enough}
\label{sec:def-why}

The Krylov--Schur machine of \cref{ch:krylovschur} already finds several
eigenpairs at once: its thick-restart loop locks each converged Ritz pair and
keeps iterating for the rest, and SLEPc's \code{EPSSolve} returns the requested
batch in one opaque call. So why a whole chapter on finding eigenpairs ``one at a
time''? Two reasons, and they pull in opposite directions until deflation
reconciles them.

The first is that \emph{locking is deflation}---the chapter is, in part, the
constructive account of a mechanism \cref{ch:krylovschur} used and named but did
not open. When Krylov--Schur ``locks'' a converged Ritz vector, it must thereafter
keep the growing Krylov basis \emph{orthogonal} to that vector, or the next
expansion step would simply re-grow the direction it just converged and the
solver would spin in place rediscovering the same mode. Locking, orthogonalizing
against the locked set, and ``purging'' a converged direction from the active
subspace are three names for one act: \emph{removing an already-found eigenvector
from the space the iteration is allowed to search}. That act is deflation, and
\cref{sec:def-locking} makes it explicit.

The second reason is sharper, and it is why deflation gets promoted from a
subroutine to a top-level algorithm. The clean black-box eigensolve of
\cref{ch:eigenvalues,ch:krylovschur} assumed a \emph{linear} or at worst
\emph{polynomial} pencil---$\Kop\vx=\lambda\Mop\vx$, or the quadratic
$(\Kop+\lambda\Cop+\lambda^2\Mop)\vx=\mathbf 0$ that linearization turns back into
a generalized linear problem (\cref{ch:eigenvalues}). For a genuinely
\emph{nonlinear} eigenproblem---one whose operator depends on $\lambda$ in a way
no polynomial captures---there is no linearization, no single library call that
returns a batch, and no Krylov basis whose Ritz values are the answers. The only
practical route known is to hunt the eigenpairs \emph{one at a time} with a
Newton-like iteration, deflating each as it converges. Deflation stops being a
bookkeeping detail inside the eigensolver and becomes the load-bearing
scaffolding of the whole method.

\begin{keyidea}
Deflation \emph{peels eigenpairs off the spectrum one at a time}: once a pair has
converged, it is projected out of the search space so the iteration is forced to
find the next one, never re-finding what it already has. For a linear pencil this
is the ``locking'' inside Krylov--Schur (\cref{ch:krylovschur}); for a
\emph{nonlinear} eigenproblem $\mat{T}(\lambda)\vx=\mathbf 0$---whose operator
\emph{depends on its own eigenvalue}---it becomes the top-level algorithm, because
no linearization and no single library call can return the modes as a batch.
\end{keyidea}

We take the two halves in order: first deflation on the familiar linear problem,
where every step can be checked against the power iteration of
\cref{ch:eigenvalues}; then the nonlinear eigenproblem and the deflated
quasi-Newton method Palace runs on it.

\section{Deflation on a linear problem}
\label{sec:def-linear}

Strip the problem to its bones. We have a symmetric matrix $\mat{A}$, we have
just found its dominant eigenpair $(\mu_1,\vu_1)$ by the power iteration of
\cref{ch:eigenvalues}, and we want the \emph{next} one, $(\mu_2,\vu_2)$. If we
simply restart power iteration from a fresh vector, it converges straight back to
$\vu_1$---the dominant direction is still dominant, and iteration has no memory of
having found it. We must change the problem so that $\vu_1$ is no longer there to
be found.

\subsection{Subtracting a rank-one piece}
\label{sec:def-ranköne}

For a symmetric matrix the cleanest construction is to \emph{subtract the found
eigenpair out of the operator}. If $\mat{A}$ is symmetric with orthonormal
eigenvectors, the dominant pair contributes a rank-one term $\mu_1\vu_1\vu_1^{\!\top}$
to its spectral decomposition $\mat{A}=\sum_i \mu_i\vu_i\vu_i^{\!\top}$. Remove
exactly that term:
\begin{equation}
  \boxed{\;\mat{A}_1 \;=\; \mat{A} \;-\; \mu_1\,\vu_1\vu_1^{\!\top}.\;}
  \label{eq:deflate-rank1}
\end{equation}
The deflated matrix $\mat{A}_1$ has \emph{exactly the same eigenvectors and
eigenvalues as $\mat{A}$, except that $\mu_1$ has been replaced by zero}. To see
it, apply $\mat{A}_1$ to each eigenvector. On $\vu_1$:
\begin{equation}
  \mat{A}_1\vu_1
    = \mat{A}\vu_1 - \mu_1\vu_1(\vu_1^{\!\top}\vu_1)
    = \mu_1\vu_1 - \mu_1\vu_1
    = \mathbf 0,
\end{equation}
using $\vu_1^{\!\top}\vu_1=1$: the dominant eigenvalue is annihilated, sent to
zero. On any other eigenvector $\vu_j$ ($j\neq1$), orthogonality
$\vu_1^{\!\top}\vu_j=0$ kills the correction term entirely:
\begin{equation}
  \mat{A}_1\vu_j
    = \mat{A}\vu_j - \mu_1\vu_1(\vu_1^{\!\top}\vu_j)
    = \mu_j\vu_j - \mathbf 0
    = \mu_j\vu_j .
\end{equation}
Every other pair survives untouched. So $\mat{A}_1$'s spectrum is
$\{0,\mu_2,\mu_3,\dots\}$: the former runner-up $\mu_2$ is now the largest in
magnitude (assuming $\abs{\mu_2}>0$), and \emph{power iteration on $\mat{A}_1$
converges to $\vu_2$}---the exact pair we wanted next. We found the second
eigenpair by finding the ``dominant'' one of a deflated operator.

\begin{figure}[t]
\centering
\begin{tikzpicture}
\begin{axis}[
    name=before,
    width=0.84\linewidth, height=30mm,
    xlabel={}, title={\footnotesize original spectrum of $\mat{A}$},
    title style={font=\footnotesize, yshift=-1.5mm},
    xmin=-0.6, xmax=9, ymin=0, ymax=1.5,
    axis lines=left, ytick=\empty, xtick={1,3,5,7.5},
    xticklabels={$\mu_4$,$\mu_3$,$\mu_2$,$\mu_1$},
    xticklabel style={font=\scriptsize}, clip=false,
  ]
  \addplot[ycomb, very thick, simBlue, mark=*, mark size=2pt] coordinates {
    (1,0.6)(3,0.85)(5,1.0)(7.5,1.25)};
  \node[font=\scriptsize\itshape, simRust, anchor=south] at (axis cs:7.5,1.25)
    {dominant: power iter.\ finds this};
\end{axis}
\begin{axis}[
    name=after, at={(before.south west)}, anchor=north west, yshift=-11mm,
    width=0.84\linewidth, height=30mm,
    title={\footnotesize after deflating $(\mu_1,\vu_1)$: spectrum of
      $\mat{A}_1=\mat{A}-\mu_1\vu_1\vu_1^{\!\top}$},
    title style={font=\footnotesize, yshift=-1.5mm},
    xmin=-0.6, xmax=9, ymin=0, ymax=1.5,
    axis lines=left, ytick=\empty, xtick={0,1,3,5},
    xticklabels={$0$,$\mu_4$,$\mu_3$,$\mu_2$},
    xticklabel style={font=\scriptsize}, clip=false,
  ]
  \addplot[ycomb, very thick, simBlue, mark=*, mark size=2pt] coordinates {
    (1,0.6)(3,0.85)(5,1.0)};
  \addplot[ycomb, very thick, simGray, mark=o, mark size=2pt] coordinates {(0,1.25)};
  \node[font=\scriptsize\itshape, simGray, anchor=south] at (axis cs:0,1.25)
    {$\mu_1\to 0$};
  \node[font=\scriptsize\itshape, simRust, anchor=south] at (axis cs:5,1.0)
    {now dominant: finds $\vu_2$};
\end{axis}
\end{tikzpicture}
\caption[Deflating one eigenpair sends it to zero]{Deflation by rank-one
subtraction. \emph{Top:} the original spectrum of $\mat{A}$; power iteration
finds the dominant pair $(\mu_1,\vu_1)$ at the right. \emph{Bottom:} subtracting
$\mu_1\vu_1\vu_1^{\!\top}$ replaces $\mu_1$ with $0$ (grey, open) and leaves every
other eigenvalue exactly where it was. The former runner-up $\mu_2$ is now the
largest in magnitude, so power iteration on the deflated operator $\mat{A}_1$
converges to $\vu_2$. Repeating---deflate, iterate, deflate---peels the spectrum
off one pair at a time.}
\label{fig:deflate-spectrum}
\end{figure}

\Cref{fig:deflate-spectrum} draws the move. Repeat it---deflate $(\mu_2,\vu_2)$ to
form $\mat{A}_2=\mat{A}_1-\mu_2\vu_2\vu_2^{\!\top}$, iterate for $\vu_3$, and so
on---and the spectrum is peeled off the top one eigenpair at a time. Each pair,
once found, is gone from the deflated operator and cannot be found again.

\subsection{Deflation as a projection: locking the search space}
\label{sec:def-locking}

The rank-one subtraction is concrete but special: it leans on symmetry (so that
the spectral decomposition has orthonormal $\vu_i$) and it modifies the operator.
The form that generalizes---and the form that matches what Krylov--Schur's locking
actually does---works on the \emph{search space} instead of the operator. Rather
than alter $\mat{A}$, we forbid the iteration from looking in the directions we
have already found.

Collect the converged eigenvectors as the columns of a tall thin matrix
$\mat{X}=[\,\vu_1\mid\cdots\mid\vu_k\,]$, the \emph{locked} set. Let $\mat{P}$ be
the orthogonal projector \emph{onto the complement} of their span:
\begin{equation}
  \mat{P} \;=\; \mat{I} - \mat{X}\mat{X}^{\!\top}
  \qquad(\text{with } \mat{X}^{\!\top}\mat{X}=\mat{I}_k).
  \label{eq:deflate-proj}
\end{equation}
Applying $\mat{P}$ to any vector strips out its component along the locked
directions: $\mat{P}\vu_i=\mathbf 0$ for each locked $\vu_i$, while $\mat{P}\vv=\vv$
for any $\vv$ already orthogonal to all of them. Now run the iteration on the
\emph{projected} operator $\mat{P}\mat{A}\mat{P}$, or---cheaper and equivalent for
iteration---apply $\mat{P}$ to the iterate after every step:
\begin{equation}
  \vv_{k+1} \;=\; \mat{P}\,\frac{\mat{A}\vv_k}{\norm{\mat{A}\vv_k}}.
  \label{eq:deflate-iterate}
\end{equation}
Each step re-orthogonalizes the iterate against the locked set, continually
sweeping out any component that has drifted back into a found direction. The
iteration physically \emph{cannot} converge to a locked eigenvector, because that
direction is projected to zero the instant it appears. It is forced to find the
next pair in the orthogonal complement.

\begin{figure}[t]
\centering
\begin{tikzpicture}[scale=1.0]
  % the locked direction (horizontal), the complement (vertical)
  \draw[->, simGray] (-0.3,0) -- (5.2,0)
    node[right,font=\scriptsize]{$\vu_1$ (locked)};
  \draw[->, simGray] (0,-0.3) -- (0,3.6)
    node[above,font=\scriptsize]{complement $\vu_1^\perp$};
  % the locked subspace shaded
  \fill[simGray!12] (0,-0.25) rectangle (5.0,0.25);
  \node[font=\scriptsize\itshape, simGray] at (3.6,-0.5) {forbidden: $\mat{P}\vu_1=\mathbf 0$};
  % an iterate that has drifted, with a component along u1
  \draw[-{Stealth[length=2.4mm]}, very thick, simBlue] (0,0) -- (2.6,2.0);
  \node[font=\scriptsize, simBlue] at (2.85,2.15) {$\mat{A}\vv_k$};
  % its projection onto the complement (drop the horizontal part)
  \draw[-{Stealth[length=2.4mm]}, very thick, simRust] (0,0) -- (0,2.0);
  \node[font=\scriptsize, simRust, anchor=east] at (-0.1,2.0) {$\mat{P}\,\mat{A}\vv_k$};
  % the projection drop (dotted)
  \draw[densely dotted, simGray] (2.6,2.0) -- (0,2.0);
  \draw[densely dotted, simGray] (2.6,2.0) -- (2.6,0);
  % the swept-out component
  \draw[-{Stealth[length=2mm]}, simGreen, thick] (0.05,0.12) -- (2.5,0.12);
  \node[font=\scriptsize, simGreen] at (1.3,0.42) {swept out each step};
  % converged next eigenvector lies in complement
  \draw[-{Stealth[length=2.4mm]}, very thick, simViolet, densely dashed] (0,0) -- (0,3.0);
  \node[font=\scriptsize, simViolet, anchor=east] at (-0.1,3.0) {$\vu_2$ (next)};
\end{tikzpicture}
\caption[Deflation projects the iterate off the locked directions]{Deflation as a
projection. The found eigenvector $\vu_1$ is \emph{locked}; the projector
$\mat{P}=\mat{I}-\mat{X}\mat{X}^{\!\top}$ kills any component along it (grey band:
$\mat{P}\vu_1=\mathbf 0$). When the operator pushes the iterate $\mat{A}\vv_k$
(blue) partly back toward the locked direction, applying $\mat{P}$ drops that
component (green) and leaves only the part in the orthogonal complement
($\mat{P}\,\mat{A}\vv_k$, rust). The iteration therefore cannot reconverge to
$\vu_1$; it is steered into the complement, where the next eigenvector $\vu_2$
(violet, dashed) lives. This is exactly the ``locking'' Krylov--Schur applies to
its converged Ritz vectors (\cref{ch:krylovschur}).}
\label{fig:deflate-project}
\end{figure}

\Cref{fig:deflate-project} shows the geometry. This projection form is what
Krylov--Schur's locking does (\cref{ch:krylovschur}): the converged Ritz vectors
are held aside, and every new Arnoldi/Lanczos column is orthogonalized against
them before it joins the active basis. ``Keep the search space orthogonal to the
locked vectors'' and ``deflate the found eigenpairs'' are the same instruction.
The block form---locking several at once into $\mat{X}$ and projecting against the
whole block---is how a solver computes eigenpairs not just one at a time but in
\emph{groups}, deflating a converged batch and continuing for the next.

\begin{notationbox}
Two equivalent ways to deflate $k$ converged eigenvectors $\mat{X}=[\vu_1\mid
\cdots\mid\vu_k]$ from a linear problem:
\begin{itemize}[nosep]
  \item \textbf{Operator deflation} (symmetric case): subtract the found spectral
  mass, $\mat{A}_k=\mat{A}-\sum_{i=1}^k\mu_i\vu_i\vu_i^{\!\top}$, sending each found
  $\mu_i$ to $0$.
  \item \textbf{Search-space deflation} (general): project with
  $\mat{P}=\mat{I}-\mat{X}\mat{X}^{\!\top}$ after each step, forbidding the
  iteration from drifting back into $\operatorname{span}(\mat{X})$.
\end{itemize}
Both enforce one contract: an already-found eigenvector cannot be found again.
The projection form is the one that survives into the nonlinear method of
\cref{sec:def-nleps}, generalized to an \emph{oblique} (non-orthogonal) projector.
\end{notationbox}

\subsection{A worked deflation}
\label{sec:def-worked1}

The whole mechanism fits in a few lines of arithmetic on a $2\times2$ matrix,
where we can carry power iteration by hand and watch the second eigenpair appear
only after deflation.

\begin{workedexample}{Find the second eigenpair by deflation}{deflate2x2}
Take the symmetric matrix
\[
  \mat{A} = \begin{psmallmatrix} 3 & 1 \\ 1 & 3 \end{psmallmatrix},
\]
whose eigenpairs we happen to know exactly: $\mu_1=4$ with
$\vu_1=\tfrac1{\sqrt2}(1,1)^{\!\top}$ and $\mu_2=2$ with
$\vu_2=\tfrac1{\sqrt2}(1,-1)^{\!\top}$. Pretend we know neither and hunt them by
power iteration.

\smallskip
\emph{Step 1: power iteration finds the dominant pair.} Start from
$\vv_0=(1,0)^{\!\top}$. One multiply gives $\mat{A}\vv_0=(3,1)^{\!\top}$; normalize
to $\vv_1=(0.949,0.316)^{\!\top}$. Again: $\mat{A}\vv_1=(3.16,2.06)^{\!\top}$,
normalize to $\vv_2=(0.838,0.546)^{\!\top}$. Again:
$\mat{A}\vv_2=(3.06,2.48)^{\!\top}\to\vv_3=(0.776,0.630)^{\!\top}$. The iterate is
swinging toward $(0.707,0.707)^{\!\top}=\vu_1$; the Rayleigh quotient
$\rho(\vv_3)=\inner{\mat{A}\vv_3}{\vv_3}=3.98$ is closing on $\mu_1=4$. Take the
converged pair $(\mu_1,\vu_1)=(4,\tfrac1{\sqrt2}(1,1)^{\!\top})$.

\smallskip
\emph{Step 2: restarting plain power iteration is useless.} If we restart from a
new random vector, it converges right back to $\vu_1$---$\mu_1=4$ is still the
dominant eigenvalue and iteration has no memory. We must deflate.

\smallskip
\emph{Step 3: deflate the found pair.} Subtract the rank-one piece
\cref{eq:deflate-rank1}, with $\vu_1\vu_1^{\!\top}=\tfrac12
\begin{psmallmatrix}1&1\\1&1\end{psmallmatrix}$:
\[
  \mat{A}_1 = \mat{A} - 4\cdot\tfrac12\begin{psmallmatrix}1&1\\1&1\end{psmallmatrix}
    = \begin{psmallmatrix}3&1\\1&3\end{psmallmatrix}
      - \begin{psmallmatrix}2&2\\2&2\end{psmallmatrix}
    = \begin{psmallmatrix}1&-1\\-1&1\end{psmallmatrix}.
\]
Check the spectrum directly: $\mat{A}_1\vu_1=\begin{psmallmatrix}1&-1\\-1&1
\end{psmallmatrix}\tfrac1{\sqrt2}(1,1)^{\!\top}=\tfrac1{\sqrt2}(0,0)^{\!\top}
=\mathbf 0$---the dominant pair is annihilated. And
$\mat{A}_1\vu_2=\begin{psmallmatrix}1&-1\\-1&1\end{psmallmatrix}\tfrac1{\sqrt2}
(1,-1)^{\!\top}=\tfrac1{\sqrt2}(2,-2)^{\!\top}=2\vu_2$---the second pair survives
with eigenvalue $2$, untouched. The deflated spectrum is $\{0,2\}$.

\smallskip
\emph{Step 4: power iteration on $\mat{A}_1$ now finds $\vu_2$.} Start again from
$\vv_0=(1,0)^{\!\top}$. One multiply: $\mat{A}_1\vv_0=(1,-1)^{\!\top}$, normalize to
$(0.707,-0.707)^{\!\top}$---which is $\vu_2$ \emph{exactly}, in a single step,
because the deflated matrix has only one nonzero eigendirection left. The Rayleigh
quotient reads $\rho=2=\mu_2$. The second eigenpair fell out the moment we
deflated the first.

The lesson generalizes past this tiny case: deflation does not make the next
eigenpair easier to \emph{compute} (the per-step work is the same), it makes the
next eigenpair the one the iteration \emph{converges to at all}.
\end{workedexample}

\begin{figure}[t]
\centering
\begin{tikzpicture}[font=\footnotesize, node distance=7mm]
  % the spectrum being peeled, lowest first
  % five spectrum stems, lowest-first (drawn explicitly, not via \foreach)
  \draw[thick, simRust] (0,0) -- (0,1.4);   \fill[simRust] (0,1.4) circle (1.6pt);
  \node[font=\scriptsize, simRust, anchor=south] at (0,1.4) {$\lambda_1$};
  \draw[thick, simRust] (1.4,0) -- (1.4,1.4); \fill[simRust] (1.4,1.4) circle (1.6pt);
  \node[font=\scriptsize, simRust, anchor=south] at (1.4,1.4) {$\lambda_2$};
  \draw[thick, simBlue] (2.8,0) -- (2.8,1.4); \fill[simBlue] (2.8,1.4) circle (1.6pt);
  \node[font=\scriptsize, simBlue, anchor=south] at (2.8,1.4) {$\lambda_3$};
  \draw[thick, simGray] (4.2,0) -- (4.2,1.4); \fill[simGray] (4.2,1.4) circle (1.6pt);
  \node[font=\scriptsize, simGray, anchor=south] at (4.2,1.4) {$\lambda_4$};
  \draw[thick, simGray] (5.6,0) -- (5.6,1.4); \fill[simGray] (5.6,1.4) circle (1.6pt);
  \node[font=\scriptsize, simGray, anchor=south] at (5.6,1.4) {$\lambda_5$};
  \draw[->, simGray] (-0.5,0) -- (6.4,0) node[right,font=\scriptsize]{$\lambda$};
  % peeling annotations
  \draw[-{Stealth[length=2mm]}, simRust, thick] (0,-0.55) -- (0,-0.05);
  \node[font=\scriptsize, simRust] at (0,-0.85) {found \#1};
  \draw[-{Stealth[length=2mm]}, simRust, thick] (1.4,-0.55) -- (1.4,-0.05);
  \node[font=\scriptsize, simRust] at (1.4,-0.85) {found \#2};
  \draw[-{Stealth[length=2mm]}, simBlue, thick] (2.8,-0.55) -- (2.8,-0.05);
  \node[font=\scriptsize, simBlue] at (2.8,-0.85) {seeking \#3};
  \node[font=\scriptsize\itshape, simGray] at (5.0,-0.85) {not yet reached};
  % the "wanted band" bracket
  \draw[decorate, decoration={brace, amplitude=4pt}, simRust]
    (-0.15,1.75) -- (1.55,1.75);
  \node[font=\scriptsize, simRust] at (0.7,2.15) {deflated \& locked};
  \draw[decorate, decoration={brace, amplitude=4pt, mirror}, simGray]
    (2.65,1.75) -- (5.75,1.75);
  \node[font=\scriptsize, simGray] at (4.2,2.15) {still in the active search};
\end{tikzpicture}
\caption[Peeling the spectrum lowest-first]{Computing many eigenpairs one at a
time. The wanted eigenvalues, lowest first, are peeled off the spectrum: pairs
$\lambda_1,\lambda_2$ are converged, deflated, and locked (rust, bracketed left);
the iteration is now hunting $\lambda_3$ (blue) in the space orthogonal to the
locked set; the higher pairs remain in the active search (grey). Each conversion
from ``seeking'' to ``found'' adds a column to the locked block $\mat{X}$ and
shrinks the search space by one direction. This lowest-first peeling is exactly
the eigenmode driver's request: the fundamental resonance and its first several
overtones, in order.}
\label{fig:peel-spectrum}
\end{figure}

\Cref{fig:peel-spectrum} shows the campaign in the large: pairs converted from
``seeking'' to ``found,'' the locked block growing by a column each time, the
spectrum peeled lowest-first. That is the picture to hold as we turn to the case
where the operator itself is moving underneath us.

\section{The nonlinear eigenproblem}
\label{sec:def-nep}

Everything so far assumed the operator was a fixed pencil. We now break that
assumption, because the physics breaks it.

\subsection{When the material depends on frequency}
\label{sec:def-dispersive}

The pencils of \cref{ch:eigenvalues} came from materials with constant
properties: a permittivity $\varepsilon$ and permeability $\mu$ that are fixed
numbers, baked once into the mass and stiffness matrices. Many real materials are
not like that. A \emph{dispersive} medium has a permittivity that depends on
frequency, $\varepsilon(\omega)$; a \emph{lossy} medium has a complex, frequency
-dependent response that absorbs energy. Plasmas, dielectric resonators with
Debye or Lorentz relaxation, conductors at microwave frequencies, metamaterials---
in all of them the material's response to the field is a function of how fast the
field oscillates.

When you discretize Maxwell's equations in such a medium, the frequency
dependence does not vanish into constant matrices. It rides along into the
operator. Writing $\lambda$ for the (squared, complex) eigenfrequency as before,
the discrete eigenproblem becomes
\begin{equation}
  \boxed{\;\mat{T}(\lambda)\,\vx \;=\; \mathbf 0,\;}
  \label{eq:nep}
\end{equation}
where the \emph{operator} $\mat{T}(\lambda)$ \emph{itself depends on the
eigenvalue $\lambda$ we are solving for}, and depends on it nonlinearly---not as
the tidy polynomial $\Kop+\lambda\Cop+\lambda^2\Mop$, but through whatever
function $\varepsilon(\lambda)$, $\nu(\lambda)$ the material model dictates. A
typical form is
\begin{equation}
  \mat{T}(\lambda) \;=\; \Kop + \lambda\Cop + \lambda^2\Mop + \mat{A}_2(\lambda),
  \label{eq:nep-form}
\end{equation}
where $\mat{A}_2(\lambda)$ is the operator-valued nonlinear part carrying the
dispersion---a rational function of $\lambda$ for a Lorentz model, a square-root
or exponential for a more exotic one. This is a \emph{nonlinear eigenproblem}, or
NEP. The word ``nonlinear'' refers to the dependence on $\lambda$, not to any
nonlinearity in $\vx$: \cref{eq:nep} is still linear in the eigenvector $\vx$ for
each fixed $\lambda$.

\begin{physicsbox}
The nonlinear eigenproblem is the mathematics of a resonance that \emph{shapes its
own cavity}. In a dispersive medium the material's response---how strongly it
polarizes, how much it absorbs---depends on the oscillation frequency. But the
oscillation frequency is precisely the resonance $\lambda$ we are trying to find.
So the operator that determines the resonance depends on the resonance itself: a
feedback loop baked into the equation. A mode ringing at one frequency sees a
different material than a mode ringing at another, and the self-consistent
frequency---the one at which the cavity, filled with the material \emph{as it
behaves at that frequency}, actually resonates---is the eigenvalue of
$\mat{T}(\lambda)\vx=\mathbf 0$. The complex part of $\lambda$ encodes the loss:
its imaginary part is the oscillation, its real part the decay, and the two
together give the resonant frequency and the quality factor $Q$
(\cref{ch:reductions}).
\end{physicsbox}

\subsection{Why linearization fails}
\label{sec:def-linfails}

\Cref{ch:eigenvalues} met the \emph{quadratic} eigenproblem
$(\Kop+\lambda\Cop+\lambda^2\Mop)\vx=\mathbf 0$ and tamed it by
\emph{linearization}: introduce $\vz=\lambda\vx$, stack $(\vx,\vz)$, and trade the
quadratic $N\times N$ problem for a linear $2N\times 2N$ generalized
eigenproblem the black-box library can solve. The trick worked because the
$\lambda$-dependence was \emph{polynomial}: a degree-two polynomial in $\lambda$
linearizes to a pencil of twice the size, a degree-$d$ polynomial to a pencil of
$d$ times the size. The companion-matrix construction is exact.

For the genuinely nonlinear $\mat{A}_2(\lambda)$ there is no such construction.
You cannot stack a finite number of auxiliary unknowns to absorb a rational or
transcendental function of $\lambda$; the ``companion'' would need infinitely many
blocks. \Cref{fig:linear-vs-nep} contrasts the two situations. A polynomial of
degree $d$ has a finite linearization (a $dN\times dN$ pencil); a true
nonlinearity has none. Approximating $\mat{A}_2(\lambda)$ by a polynomial and
linearizing \emph{that} is a real technique (rational and polynomial
approximation methods underlie much of modern NEP solving), but it trades the
exact problem for an approximate one and inflates the dimension. Palace takes the
other road: confront the nonlinearity directly with a Newton iteration on
$\lambda$, and use deflation to get more than one eigenpair.

\begin{figure}[t]
\centering
\begin{tikzpicture}[font=\footnotesize]
  % ===== LEFT: linear/polynomial pencil — fixed operator =====
  \begin{scope}
    \node[font=\small\bfseries, simBlue] at (1.4,2.9) {linear / polynomial pencil};
    \node[stage, minimum width=30mm, minimum height=14mm] (op) at (1.4,1.4)
      {$\Kop+\lambda\Cop+\lambda^2\Mop$};
    \node[font=\scriptsize\itshape, simGray] at (1.4,0.55)
      {operator built from \emph{fixed} matrices};
    \node[data, minimum width=20mm] (lam) at (1.4,-0.4) {$\lambda$};
    % lambda only scales the fixed blocks
    \draw[flow, simGray] (lam) -- node[right,font=\scriptsize]{scales blocks} (op);
    \node[font=\scriptsize, simGreen, align=center] at (1.4,-1.4)
      {\emph{linearizes:} $2N\!\times\!2N$\\ pencil, one library call};
  \end{scope}
  % divider
  \draw[simGray!50, densely dashed] (3.6,-1.8) -- (3.6,3.1);
  % ===== RIGHT: nonlinear T(lambda) — operator depends on its own eigenvalue =====
  \begin{scope}[xshift=4.6cm]
    \node[font=\small\bfseries, simRust] at (1.7,2.9) {nonlinear $\mat{T}(\lambda)$};
    \node[stage, minimum width=34mm, minimum height=14mm, draw=simRust, fill=simBgRust]
      (opn) at (1.7,1.4) {$\mat{T}(\lambda)=\cdots+\mat{A}_2(\lambda)$};
    \node[data, minimum width=20mm] (lamn) at (1.7,-0.6) {$\lambda$};
    % feedback: lambda feeds the operator, the operator's eigenvalue IS lambda
    \draw[flow, simRust] (lamn.east) to[out=0,in=-60]
      node[right,font=\scriptsize,align=left,pos=0.4]{enters the\\operator} (opn.south east);
    \draw[backflow, simRust] (opn.south west) to[out=-120,in=180]
      node[left,font=\scriptsize,align=right,pos=0.6]{whose eigenvalue\\\emph{is} $\lambda$} (lamn.west);
    \node[font=\scriptsize, simRust, align=center] at (1.7,-1.7)
      {\emph{no linearization:} solve\\ $\mat{T}(\lambda)\vx=\mathbf 0$ directly};
  \end{scope}
\end{tikzpicture}
\caption[Linear pencil versus nonlinear $\mat{T}(\lambda)$]{The structural break.
\emph{Left:} a linear or polynomial pencil is built from \emph{fixed} matrices
$\Kop,\Cop,\Mop$; the eigenvalue $\lambda$ only \emph{scales} those blocks, so the
problem linearizes exactly to a larger pencil and the black-box library of
\cref{ch:krylovschur} solves it in one call. \emph{Right:} for a dispersive
material the operator $\mat{T}(\lambda)$ \emph{depends on $\lambda$ itself}
through the nonlinear part $\mat{A}_2(\lambda)$---a feedback loop, since the
$\lambda$ feeding the operator is the very eigenvalue we seek. No finite stacking
of auxiliary unknowns can absorb a non-polynomial dependence, so linearization
fails and we must attack $\mat{T}(\lambda)\vx=\mathbf 0$ directly.}
\label{fig:linear-vs-nep}
\end{figure}

\section{Deflated quasi-Newton for the NEP}
\label{sec:def-nleps}

Palace solves the nonlinear eigenproblem with a \emph{deflated quasi-Newton}
method---the scheme SLEPc's NEP solver uses, with minimality index one, following
Effenberger~\citep{effenberger2013}. The idea is exactly the two halves of this
chapter, joined: \emph{Newton's method to converge one eigenpair} of the
nonlinear system, and \emph{deflation to lock it and move to the next}. We build
it in that order.

\subsection{Newton on the eigenvalue}
\label{sec:def-newton}

Forget deflation for a moment and chase a single eigenpair of
$\mat{T}(\lambda)\vx=\mathbf 0$. The unknowns are the pair $(\vx,\lambda)$---an
$N$-vector and a scalar---and \cref{eq:nep} is $N$ equations. That is one equation
short of a square system: any scalar multiple of an eigenvector is also an
eigenvector, so $\vx$ is determined only up to scale, and we need a
\emph{normalization} to pin it down. Add one: fix the size and phase of $\vx$ by
requiring $\vc^{\!\top}\vx=1$ for some fixed vector $\vc$. Now we have $N+1$
equations in $N+1$ unknowns,
\begin{equation}
  \mat{F}(\vx,\lambda) \;=\;
  \begin{pmatrix} \mat{T}(\lambda)\,\vx \\[2pt] \vc^{\!\top}\vx - 1 \end{pmatrix}
  \;=\; \mathbf 0,
  \label{eq:nep-system}
\end{equation}
a square \emph{nonlinear system} we can hand to Newton's method---the same
quasi-Newton machinery that drove the nonlinear solves of \cref{ch:cg}, now with
the eigenvalue itself among the unknowns.

Newton's method linearizes $\mat{F}$ about the current guess and solves for the
correction. The Jacobian of \cref{eq:nep-system} with respect to
$(\vx,\lambda)$---differentiating each block---is the bordered matrix
\begin{equation}
  \mat{J}(\vx,\lambda) =
  \begin{pmatrix}
    \mat{T}(\lambda) & \mat{T}'(\lambda)\,\vx \\[2pt]
    \vc^{\!\top} & 0
  \end{pmatrix},
  \label{eq:nep-jacobian}
\end{equation}
where $\mat{T}'(\lambda)=\mathrm{d}\mat{T}/\mathrm{d}\lambda$ is the derivative of
the operator with respect to the eigenvalue---the term that carries the
nonlinearity into the Jacobian. The Newton step solves
$\mat{J}(\vx,\lambda)\,(\Delta\vx,\Delta\lambda)^{\!\top}=-\mat{F}(\vx,\lambda)$
and updates $(\vx,\lambda)\leftarrow(\vx,\lambda)+(\Delta\vx,\Delta\lambda)$. The
\emph{quasi}-Newton variant Palace uses replaces the exact Jacobian action with a
cheaper lagged approximation (and an inexact inner tolerance), in the same spirit
as the quasi-Newton methods of \cref{ch:cg}; the structure---residual, Jacobian
solve, update---is identical.

\begin{keyidea}
The nonlinear eigenproblem $\mat{T}(\lambda)\vx=\mathbf 0$ is turned into a square
\emph{nonlinear system} in the unknowns $(\vx,\lambda)$ by appending a
normalization $\vc^{\!\top}\vx=1$, and solved by \emph{Newton's method---with the
eigenvalue $\lambda$ itself as one of the unknowns Newton drives}. Each Newton
step evaluates a residual, solves a bordered linear system built from
$\mat{T}(\lambda)$ and its derivative $\mat{T}'(\lambda)$, and updates both the
eigenvector and the eigenvalue. ``Newton on the eigenvalue'' is the engine; we
deflate around it to get more than one mode.
\end{keyidea}

The next worked example makes the eigenvalue-Newton concrete on a scalar problem
small enough to compute by hand, where ``Newton on $\lambda$'' is visibly just
Newton's method on a scalar equation.

\begin{workedexample}{Newton on a scalar nonlinear eigenproblem}{nepscalar}
Strip the NEP to one dimension, so $\vx$ is a scalar $x$ and $\mat{T}(\lambda)$ is
a scalar function $T(\lambda)$. Then $\mat{T}(\lambda)x=0$ with $x\neq0$ forces
$T(\lambda)=0$: \emph{the nonlinear eigenproblem is just a root-find for
$\lambda$}, and Newton on the eigenvalue is plain scalar Newton. Take the
nonlinearity
\[
  T(\lambda) \;=\; \lambda^2 - g(\lambda),
  \qquad g(\lambda) = 2 + \tfrac{1}{2}\cos\lambda,
\]
a quadratic with a small dispersive wobble $g(\lambda)$ standing in for the
material's frequency dependence. (Were $g$ constant, this would be an ordinary
quadratic eigenvalue $\lambda=\sqrt g$; the $\cos\lambda$ is what makes it
genuinely nonlinear.) We seek the smallest positive root.

\smallskip
\emph{The Newton step.} With $T'(\lambda)=2\lambda+\tfrac12\sin\lambda$, the scalar
Newton update is
\[
  \lambda_{k+1} = \lambda_k - \frac{T(\lambda_k)}{T'(\lambda_k)}
    = \lambda_k - \frac{\lambda_k^2 - 2 - \tfrac12\cos\lambda_k}
                       {2\lambda_k + \tfrac12\sin\lambda_k}.
\]
This is exactly the bordered Newton step \cref{eq:nep-jacobian} in one dimension:
the ``Jacobian'' is the scalar $T'(\lambda)$, the residual is $T(\lambda)$, and
the normalization holds $x=1$ so $\Delta x=0$ and only $\lambda$ moves.

\smallskip
\emph{Iterate from $\lambda_0=1.5$.}
\[
  \begin{array}{c|ccc}
    k & \lambda_k & T(\lambda_k) & \lambda_{k+1} \\\hline
    0 & 1.5000 & 0.2147 - 0\!=\!-0.2147 & 1.5713 \\
    1 & 1.5713 & 0.0205 & 1.5651 \\
    2 & 1.5651 & 0.00018 & 1.56504 \\
    3 & 1.56504 & 1.4\!\times\!10^{-8} & 1.565036
  \end{array}
\]
(Here $T(1.5)=2.25-2-\tfrac12\cos1.5=0.25-0.0354=0.2146$, and the step
$1.5-0.2146/3.0046=1.5714$, matching the table to rounding.) The residual
$T(\lambda_k)$ is driven to zero \emph{quadratically}---roughly doubling its
correct digits each step, the signature of Newton convergence---and settles at the
self-consistent eigenvalue $\lambda^\star\approx1.5650$, the frequency at which
$\lambda^2$ equals the material response $g(\lambda)$ \emph{evaluated at that very
frequency}. The fixed point of the feedback loop of \cref{fig:linear-vs-nep} is a
root, and Newton finds it.

\smallskip
\emph{Reading the residual figure.} \Cref{fig:newton-residual} plots
$\abs{T(\lambda_k)}$ against the step $k$ on a log axis: a straight line
steepening downward, the quadratic plunge. In the full $N$-dimensional problem the
scalar division becomes the bordered linear solve \cref{eq:nep-jacobian}, but the
convergence story is the one on this page.
\end{workedexample}

\begin{figure}[t]
\centering
\begin{tikzpicture}
\begin{axis}[
  width=11cm, height=5.6cm,
  xlabel={quasi-Newton step $k$},
  ylabel={residual $\abs{T(\lambda_k)}$},
  ymode=log, ymin=1e-10, ymax=1,
  xmin=0, xmax=4,
  grid=both, grid style={simGray!18},
  xtick={0,1,2,3,4},
  label style={font=\footnotesize}, tick label style={font=\scriptsize},
  legend pos=north east, legend style={font=\scriptsize, fill=white, draw=simGray!40},
  legend cell align=left,
]
% the residual from the worked example: quadratic convergence
\addplot[simRust, very thick, mark=*, mark size=2pt] coordinates {
  (0,0.2147)(1,0.0205)(2,1.8e-4)(3,1.4e-8)};
\addlegendentry{$\abs{T(\lambda_k)}$ (Newton on $\lambda$)}
% reference: linear convergence (a shallower line) for contrast
\addplot[simBlue, very thick, densely dashed, mark=square, mark size=1.5pt] coordinates {
  (0,0.2147)(1,0.064)(2,0.019)(3,0.0058)(4,0.0017)};
\addlegendentry{linear convergence (for contrast)}
% tolerance line
\addplot[simGreen, densely dotted, thick, domain=0:4, samples=2] {1e-8};
\node[font=\scriptsize, simGreen, anchor=west] at (axis cs:0.1,3e-8) {tolerance};
\end{axis}
\end{tikzpicture}
\caption[Quasi-Newton drives the eigenvalue residual to zero]{Newton on the
eigenvalue, convergence history from \cref{wex:nepscalar}. The residual
$\abs{T(\lambda_k)}$ (rust) falls \emph{quadratically}---each step roughly squares
the previous error, so the count of correct digits doubles per step and the
residual plunges through the tolerance (green dotted) in three or four steps. The
dashed blue line shows merely \emph{linear} convergence for contrast: the
straight-line decay of a fixed-rate iteration, far slower. This quadratic plunge
is the payoff of solving the bordered Newton system \cref{eq:nep-jacobian} each
step rather than fixed-point iterating. In the full problem each rust point costs
one bordered linear solve; in the scalar case it costs one division.}
\label{fig:newton-residual}
\end{figure}

\subsection{Deflating through an extended problem}
\label{sec:def-extended}

Newton converges \emph{one} eigenpair. To get the next, we deflate---but the
rank-one operator subtraction of \cref{sec:def-ranköne} relied on symmetry and a
fixed operator, neither of which a nonlinear, non-symmetric $\mat{T}(\lambda)$
offers. The NEP deflation is instead built into the \emph{system Newton solves},
by augmenting it.

Suppose $k$ eigenpairs have already converged, collected as the columns of
$\mat{X}\in\Complex^{N\times k}$ (the locked invariant basis) together with a
small $k\times k$ matrix $\mat{H}$ recording how the operator acts \emph{within}
that found subspace (its Rayleigh-quotient block). Effenberger's construction
forms an \emph{extended} nonlinear eigenproblem of size $n+k$---the original $N$
unknowns of $\vx$ plus $k$ extra coordinates $\vx_2\in\Complex^k$ that measure the
new pair's overlap with the locked block:
\begin{equation}
  \widetilde{\mat{T}}(\lambda)
  \begin{pmatrix}\vx\\\vx_2\end{pmatrix}
  =\mathbf 0,
  \qquad
  \widetilde{\mat{T}}(\lambda) =
  \begin{pmatrix}
    \mat{T}(\lambda) & \mat{U}(\lambda) \\
    \mat{X}^{\!\mathsf H} & \mathbf 0
  \end{pmatrix},
  \label{eq:nep-extended}
\end{equation}
where the coupling block $\mat{U}(\lambda)=\mat{T}(\lambda)\,\mat{X}\,
(\lambda\mat{I}-\mat{H})^{-1}$ ties the new unknown to the locked pairs through a
\emph{Schur-modified} factor
\begin{equation}
  \boxed{\;\mat{S} \;=\; \lambda\mat{I} - \mat{H}\;}
  \label{eq:nep-schur}
\end{equation}
built from the eigenvalues of the already-found pairs (encoded in $\mat{H}$). The
construction has a clean property, the whole reason for the augmentation: the
extended problem \cref{eq:nep-extended} has \emph{exactly the eigenpairs of the
original $\mat{T}(\lambda)$ that are not already in $\mat{X}$}. The locked pairs
have been removed from the extended spectrum---deflated---so Newton on the extended
system cannot reconverge to them and is driven to a genuinely new pair. The bottom
block-row $\mat{X}^{\!\mathsf H}\vx=\mathbf 0$ is the familiar
``orthogonal-to-the-locked-set'' constraint of \cref{sec:def-locking}, now an
explicit equation in the augmented system rather than a projection applied by
hand; the Schur factor $\mat{S}=\lambda\mat{I}-\mat{H}$ is what makes the
deflation respect the nonlinear, $\lambda$-dependent structure.

\begin{figure}[t]
\centering
\begin{tikzpicture}[font=\footnotesize]
  % the extended (n+k) x (n+k) block matrix
  % big block T(lambda): n x n
  \draw[thick, fill=simBgBlue, draw=simBlue] (0,0) rectangle (3.0,3.0);
  \node[font=\small\bfseries, simBlue] at (1.5,1.6) {$\mat{T}(\lambda)$};
  \node[font=\scriptsize, simGray] at (1.5,1.0) {$n\times n$, sparse, huge};
  % coupling block U(lambda): n x k  (right strip)
  \draw[thick, fill=simBgRust, draw=simRust] (3.0,0) rectangle (3.9,3.0);
  \node[font=\scriptsize, simRust, rotate=90] at (3.45,1.5) {$\mat{U}(\lambda)$};
  % locked block X^H: k x n  (bottom strip)
  \draw[thick, fill=simBgGreen, draw=simGreen] (0,-0.9) rectangle (3.0,0.0);
  \node[font=\scriptsize, simGreen] at (1.5,-0.45) {$\mat{X}^{\!\mathsf H}$ \ (locked, $k\times n$)};
  % small dense block: k x k
  \draw[thick, fill=simBgGold, draw=simGold] (3.0,-0.9) rectangle (3.9,0.0);
  \node[font=\scriptsize, simGold] at (3.45,-0.45) {$\mathbf 0$};
  % labels / braces
  \draw[decorate, decoration={brace, amplitude=4pt}, simBlue] (-0.1,3.05) -- (2.95,3.05);
  \node[font=\scriptsize, simBlue] at (1.45,3.45) {original $n$ unknowns $\vx$};
  \draw[decorate, decoration={brace, amplitude=4pt}, simRust] (3.05,3.05) -- (3.9,3.05);
  \node[font=\scriptsize, simRust, anchor=west] at (4.05,3.05) {$k$ coords $\vx_2$};
  % the Schur factor callout
  \node[font=\scriptsize, simInk, align=left, anchor=west] at (4.6,1.6)
    {$\mat{U}(\lambda)=\mat{T}(\lambda)\,\mat{X}\,\mat{S}^{-1}$,\\[2pt]
     Schur factor $\mat{S}=\lambda\mat{I}-\mat{H}$,\\[2pt]
     $\mat{H}$: locked eigenvalues};
  % the n+k size label
  \draw[decorate, decoration={brace, amplitude=5pt, mirror}, simGray] (-0.25,-0.95) -- (-0.25,3.0);
  \node[font=\scriptsize, simGray, rotate=90, anchor=south] at (-0.65,1.05) {size $n+k$};
\end{tikzpicture}
\caption[The deflated extended problem of size $n+k$]{Deflation by augmentation
(Effenberger). To deflate the $k$ already-found pairs $\mat{X}$, the original
$n\times n$ nonlinear operator $\mat{T}(\lambda)$ (blue) is bordered into an
\emph{extended} problem of size $n+k$: a coupling column $\mat{U}(\lambda)$
(rust), the locked block $\mat{X}^{\!\mathsf H}$ (green, the
``orthogonal-to-the-found-set'' constraint of \cref{sec:def-locking} made an
explicit equation), and a small $k\times k$ corner. The coupling
$\mat{U}(\lambda)=\mat{T}(\lambda)\mat{X}\mat{S}^{-1}$ uses the
\emph{Schur-modified} factor $\mat{S}=\lambda\mat{I}-\mat{H}$ built from the
locked eigenvalues in $\mat{H}$, so the deflation respects the nonlinear
$\lambda$-dependence. The extended spectrum is exactly the original spectrum
\emph{minus} the locked pairs: Newton on the augmented system finds a genuinely
new eigenpair, never a re-found one. Each converged pair grows $\mat{X}$ by one
column and $\mat{H}$ by one row and column---$k\to k+1$.}
\label{fig:extended-problem}
\end{figure}

\Cref{fig:extended-problem} draws the augmented operator. The bookkeeping that
makes it an algorithm is the growth of the locked block: each time the inner
Newton iteration converges a pair $(\lambda,\vx)$, the eigenvector is normalized
and appended as a new column of $\mat{X}$, the matrix $\mat{H}$ is \emph{bordered}
by one new row and column recording the new pair, and $k$ increments. The next
inner solve runs against the now-larger extended problem, which has one more pair
deflated out. ``Minimality index one'' is the technical guarantee that this
augmentation stays as small as possible---it adds exactly $k$ extra coordinates
for $k$ found pairs, never more---so the extended problem never bloats. The locked
pairs are invariant: the outer loop only \emph{appends} to $\mat{X}$ and
\emph{borders} $\mat{H}$, never disturbing what is already there, which is the
structural reason the one-at-a-time peeling of \cref{fig:peel-spectrum} is
correct.

\begin{figure}[t]
\centering
\begin{tikzpicture}[font=\footnotesize, node distance=8mm]
  % outer loop: one-at-a-time over eigenpairs
  \node[stage, minimum width=26mm, minimum height=11mm] (seed)
    {\textbf{seed}\\[1pt]\scriptsize linear-eigensolve\\\scriptsize guess
     (\cref{ch:krylovschur})};
  \node[stage, minimum width=30mm, minimum height=13mm, right=14mm of seed] (newton)
    {\textbf{inner quasi-Newton}\\[1pt]\scriptsize residual $\rhd$ Jacobian solve\\
     \scriptsize $\rhd$ $\lambda$-correction $\rhd$ backtrack};
  \node[stage, minimum width=24mm, minimum height=11mm, right=14mm of newton] (deflate)
    {\textbf{deflate}\\[1pt]\scriptsize append col to $\mat{X}$,\\
     \scriptsize border $\mat{H}$, $k\!+\!1$};
  \node[product, right=12mm of deflate] (done) {all $nev$\\found?};
  \draw[flow] (seed) -- (newton);
  \draw[flow] (newton) -- node[above,font=\scriptsize]{converged} (deflate);
  \draw[flow] (deflate) -- (done);
  % the inner Newton self-loop
  \draw[backflow] (newton.north) to[out=120,in=60,looseness=4]
    node[above,font=\scriptsize]{not yet} (newton.north);
  % the outer deflation loop back to seed
  \draw[backflow] (done.south) |- node[pos=0.25,below,font=\scriptsize,simRust]
    {next pair, now deflated against the grown $\mat{X}$} (seed.south);
  % the extern callout: inner linear solve
  \node[extern, below=13mm of newton, align=center] (solve)
    {\scriptsize each Newton step: one\\\scriptsize bordered, deflated linear solve\\
     \scriptsize (\cref{fig:dense-sparse-split})};
  \draw[refedge] (solve.north) -- (newton.south);
\end{tikzpicture}
\caption[The deflated quasi-Newton loop, nested]{The full method as two nested
loops. The \emph{outer} loop peels eigenpairs one at a time
(\cref{fig:peel-spectrum}): \textbf{seed} a guess from the linear eigensolver of
\cref{ch:krylovschur}, run the \textbf{inner quasi-Newton} iteration to converge
one pair against the current deflation block $\mat{X}$, then \textbf{deflate}---
append the new eigenvector to $\mat{X}$, border $\mat{H}$, increment $k$---and loop
back for the next pair, now solving the extended problem with one more pair
removed. The inner loop is Newton on the eigenvalue (\cref{sec:def-newton}):
each step evaluates the deflated residual, solves the bordered Jacobian system
(dashed callout, opened in \cref{fig:dense-sparse-split}), corrects
$(\vx,\lambda)$, and backtracks for stability, repeating until the residual falls
below tolerance. The loop terminates when $nev$ pairs are locked. This is the
\code{QuasiNewtonSolver} Palace runs.}
\label{fig:nleps-loop}
\end{figure}

\Cref{fig:nleps-loop} assembles the whole method: an outer one-at-a-time
deflation loop wrapping an inner quasi-Newton iteration, with the deflation block
growing by a column between outer iterations. The outer loop is seeded from the
linear eigensolver of \cref{ch:krylovschur}---a cheap linear-pencil approximation
gives Newton its starting guess, which the inner iteration then corrects to a true
nonlinear eigenpair (the seed need not be accurate; Newton converges it).

\section{The inner solve: dense-small meets sparse-large}
\label{sec:def-innersolve}

Each quasi-Newton step solves the bordered, deflated linear system---the action of
the extended Jacobian \cref{eq:nep-jacobian} on the augmented unknowns. This
system has a structure worth seeing clearly, because it is where two of the
linear-algebra tools of this Part meet on the same step, each handling the part it
is built for.

The bordered system splits along the $n+k$ partition of
\cref{fig:extended-problem}. The \emph{large} block is the $n\times n$ action of
$\mat{T}(\lambda)$---sparse, with millions of rows, exactly the kind of operator
\cref{ch:krylov,ch:gmres} taught us to solve iteratively. The \emph{small} block
is the $k\times k$ Schur/Gram coordinate system---dense, but tiny, with $k$ at
most a few hundred even after many modes are found. These call for two completely
different solvers:
\begin{itemize}[nosep]
  \item the large sparse block is inverted by an iterative \code{ksp\_solve}
  (\cref{ch:krylov})---a preconditioned Krylov method, never a factorization,
  because the matrix is far too large to factor;
  \item the small dense block---the $k\times k$ Gram matrix $\mat{X}^{\!\mathsf H}
  \mat{X}$ and the Schur factor $\mat{S}=\lambda\mat{I}-\mat{H}$---is inverted by a
  direct \code{lu\_solve}: a dense LU factorization, which on a $k\times k$ matrix
  with $k\sim10$--$100$ is utterly negligible.
\end{itemize}
The block elimination that couples them is the standard Schur-complement reduction
of a bordered system (\cref{ch:gmres} met the same shape in the bordered GMRES
least-squares step): solve the big block, condense its contribution into the small
block, solve the small dense block for the $k$ coordinates, and back-substitute.
The expensive part---the one sparse solve---is exactly one \code{ksp\_solve}; the
deflation costs only the trivial dense work on top.

\begin{figure}[t]
\centering
\begin{tikzpicture}[font=\footnotesize, node distance=9mm]
  % the bordered system, drawn as two coupled solves
  \node[stage, minimum width=40mm, minimum height=22mm, draw=simBlue, fill=simBgBlue]
    (big) at (0,0) {};
  \node[font=\small\bfseries, simBlue] at (big.north) [yshift=-4mm] {large sparse block};
  \node[font=\scriptsize, simInk, align=center] at (big.center) [yshift=1mm]
    {$\mat{T}(\lambda)$, \ $n\times n$\\[1pt]$n\sim 10^6$, sparse};
  \node[extern, minimum width=30mm, below=2mm of big.south, anchor=north]
    (kspsolve) {\code{ksp\_solve} (\cref{ch:krylov})\\[1pt]\scriptsize preconditioned Krylov,
      \emph{no} factorization};
  % the small dense block
  \node[stage, minimum width=26mm, minimum height=16mm, draw=simGold, fill=simBgGold,
    right=24mm of big] (small) {};
  \node[font=\small\bfseries, simGold] at (small.north) [yshift=-4mm] {small dense block};
  \node[font=\scriptsize, simInk, align=center] at (small.center) [yshift=1mm]
    {$\mat{S}=\lambda\mat{I}-\mat{H}$,\\[1pt]$\mat{X}^{\!\mathsf H}\mat{X}$, \ $k\times k$\\[1pt]
     $k\sim 10$--$100$, dense};
  \node[opbox, minimum width=24mm, below=2mm of small.south, anchor=north,
    draw=simGold] (lusolve) {\code{lu\_solve}\\[1pt]\scriptsize dense LU, negligible};
  % the Schur-complement coupling
  \draw[flow] (big.east) -- node[above,font=\scriptsize,align=center]
    {condense\\(Schur compl.)} (small.west);
  \draw[backflow] (small.south west) to[out=-150,in=-30]
    node[below,font=\scriptsize]{back-substitute the $k$ coords} (big.south east);
  % the cost annotation
  \node[font=\scriptsize\itshape, simRust, align=center] at (0,-3.4)
    {\textbf{the cost lives here:} one sparse solve/step};
  \node[font=\scriptsize\itshape, simGreen, align=center] at (5.0,-3.4)
    {\textbf{essentially free:}\\ a $k\times k$ factorization};
\end{tikzpicture}
\caption[The dense-small / sparse-large solve split]{One quasi-Newton step's
bordered linear solve, split by block. The deflated extended Jacobian
\cref{eq:nep-jacobian} partitions into a \emph{large sparse} block---the
$n\times n$ action of $\mat{T}(\lambda)$, with millions of unknowns, inverted by a
preconditioned iterative \code{ksp\_solve} (\cref{ch:krylov}; never factored)---and
a \emph{small dense} block---the $k\times k$ Schur factor
$\mat{S}=\lambda\mat{I}-\mat{H}$ and Gram matrix $\mat{X}^{\!\mathsf H}\mat{X}$,
inverted by a direct dense \code{lu\_solve}. The Schur-complement reduction
condenses the big block into the small one, solves the $k$ dense coordinates, and
back-substitutes. The entire deflation overhead is the negligible $k\times k$
factorization; the cost of the step is the single sparse solve it would have cost
without any deflation at all. Big-and-sparse iterative, small-and-dense direct:
the recurring division of labour of this Part.}
\label{fig:dense-sparse-split}
\end{figure}

\Cref{fig:dense-sparse-split} draws the division of labour. It is the same split
this Part keeps returning to---iterative Krylov for the huge sparse operator,
direct dense factorization for the tiny coordinate system---and the deflated
quasi-Newton step is where the two sit side by side in a single bordered solve, the
sparse \code{ksp\_solve} carrying essentially all the cost and the dense
\code{lu\_solve} carrying the deflation bookkeeping for free.

\section{Where this lands in Palace}
\label{sec:def-palace}

Palace's eigenmode driver is, by default, the clean linear path of
\cref{ch:eigenvalues,ch:krylovschur}: assemble $(\Kop,\Mop)$, choose a shift,
hand the pencil to SLEPc or ARPACK, read back the modes. The deflated quasi-Newton
method of this chapter is the \emph{nonlinear} branch of that same driver, taken
when the configuration declares frequency-dependent (lossy, dispersive)
materials---the case where the operator $\mat{T}(\lambda)$ genuinely depends on
$\lambda$ and no linearization exists.

In the source, this branch is a Palace-authored solver,
\code{QuasiNewtonSolver}, that drives the SLEPc-NEP-style deflation scheme
(minimality index one, Effenberger~\citep{effenberger2013}; the disguised
quasi-Newton variant of Jarlebring--Koskela--Mele). It is one of the few places
where the eigensolver's outer loop is \emph{Palace's} code rather than the
library's: where the linear eigensolve hands its whole iteration to an opaque
\code{EPSSolve} (\cref{ch:eigenvalues}), the nonlinear path runs the outer
deflation loop in Palace, calling the library only for the inner linear solves and
the linear-eigensolve seed. The structure of \cref{fig:nleps-loop}---outer
one-at-a-time deflation, inner quasi-Newton, the dense/sparse bordered solve---is
that source rendered as a composition, the outer sweep an explicit value-threaded
recursion over the growing locked block.

The contrast with the linear path is the contrast that organizes this whole Part.
The linear eigenmode solve is the fourth \emph{solve corner} of
\cref{ch:situating}: an opaque black-box call, no Palace loop. The nonlinear
eigenmode solve is a Palace-authored \emph{fold}---the same shape as the transient
time-march---that happens to wrap a black-box linear solve and a black-box linear
eigensolve inside each step. It is the deepest composition the eigenmode analysis
reaches: an outer Palace deflation loop, an inner Palace quasi-Newton loop, an
inner-inner sparse \code{ksp\_solve}, all to find the resonances of a cavity whose
material fights back at every frequency.

\subsection{Reading the results}
\label{sec:def-results}

The output is the same shape as the linear eigenmode analysis, read the same way
(\cref{ch:eigenvalues,ch:reductions}). Each converged pair $(\lambda_i,\vx_i)$
carries a \emph{complex} eigenvalue---necessarily complex, because the dispersive
material is lossy. Its imaginary part gives the oscillation frequency and its real
part the decay rate, and the two together give the resonant frequency
$f_i=\Im\lambda_i/2\pi$ (up to convention) and the quality factor
\begin{equation}
  Q_i \;=\; \frac{\abs{\Im\lambda_i}}{2\,\abs{\Re\lambda_i}},
\end{equation}
exactly as the lossy quadratic case of \cref{ch:eigenvalues} did. The eigenvector
$\vx_i$ is the mode field. The reduction from complex eigenpair to
frequency-and-$Q$ is the output postprocessing of \cref{ch:reductions}; the
nonlinear path is one more producer feeding it. So the deflated quasi-Newton
method, for all its internal machinery, hands the engineer the same table the
black-box eigensolve does---resonant frequencies and quality factors---computed
for the materials the black-box path could not represent.

\begin{pitfall}
The nonlinear path's eigenvalues are \emph{always} complex and must be read as
such; do not discard the real part as numerical noise the way one might for a
nearly-lossless linear mode. In the NEP the real part of $\lambda$ \emph{is the
loss}---it is the physical content the dispersive material was added to capture,
and it is what distinguishes a sharp high-$Q$ resonance from a broad lossy one.
A second trap is the seed: the linear-eigensolve guess that starts each outer
iteration is computed from a \emph{constant-material} approximation, so it can sit
far from the true nonlinear eigenvalue. That is expected---Newton's quadratic
convergence (\cref{fig:newton-residual}) closes the gap---but it means the seed's
eigenvalue should never be reported as the answer; only the converged
post-Newton pair is physical.
\end{pitfall}

\summarysection
The engineer wants many resonances, not one, and \emph{deflation} is how an
eigensolver delivers them: once a pair has converged it is \emph{locked}---projected
out of the search space so the iteration is forced to find the next one, peeling
the spectrum off lowest-first. On a symmetric linear problem deflation is a
rank-one operator subtraction $\mat{A}_1=\mat{A}-\mu_1\vu_1\vu_1^{\!\top}$ that
sends the found eigenvalue to zero; in general it is a projection
$\mat{P}=\mat{I}-\mat{X}\mat{X}^{\!\top}$ against the locked block---the same
``locking'' Krylov--Schur applies to its converged Ritz vectors. When the material
is dispersive the operator depends on its own eigenvalue,
$\mat{T}(\lambda)\vx=\mathbf 0$, a \emph{nonlinear eigenproblem} that no
linearization can reduce to a pencil. Palace solves it by \emph{deflated
quasi-Newton}: append a normalization to make $(\vx,\lambda)$ a square nonlinear
system, run \emph{Newton on the eigenvalue}, and deflate each found pair through an
\emph{extended problem} of size $n+k$ with the Schur-modified factor
$\mat{S}=\lambda\mat{I}-\mat{H}$. Each Newton step solves a bordered system that
splits into a large sparse block (an iterative \code{ksp\_solve}) and a tiny dense
block (a direct \code{lu\_solve}). The result is a set of complex eigenpairs read,
like the linear path, as resonant frequencies and quality factors.

\furtherreading
The deflated quasi-Newton method and its minimality-preserving extended-problem
construction are due to Effenberger~\citep{effenberger2013}, whose paper is the
direct source for the SLEPc-NEP scheme Palace runs; it is the place to go for the
proof that the extended problem \cref{eq:nep-extended} has exactly the
non-locked eigenpairs and for the minimality-index machinery only sketched here.
Saad's eigenvalue text~\citep{saad2011eig} is the standard reference for deflation
on linear problems---the rank-one subtraction, the projection form, and their
numerical conditioning---and for the broader landscape of nonlinear eigenvalue
methods (contour integration, rational approximation) that Palace's Newton route
is one branch of. Stewart~\citep{stewart2001krylovschur} connects the locking of
the previous chapter to the deflation of this one: thick-restart Krylov--Schur's
treatment of converged Ritz vectors is deflation on the linear pencil, and reading
his account alongside \cref{sec:def-locking} shows the two chapters describing one
mechanism from two sides.

\exercisessection
\begin{enumerate}[nosep]
  \item \textbf{Deflation by hand.} For
  $\mat{A}=\begin{psmallmatrix}5&2\\2&5\end{psmallmatrix}$, find the dominant
  eigenpair $(\mu_1,\vu_1)$ by inspection (the eigenvectors are
  $(1,1)^{\!\top}/\sqrt2$ and $(1,-1)^{\!\top}/\sqrt2$). Form the deflated matrix
  $\mat{A}_1=\mat{A}-\mu_1\vu_1\vu_1^{\!\top}$ explicitly and verify its spectrum is
  $\{0,\mu_2\}$. What does one step of power iteration on $\mat{A}_1$ from
  $(1,0)^{\!\top}$ return?
  \item \textbf{Projection equals subtraction.} Show that for a symmetric
  $\mat{A}$ with a single locked unit eigenvector $\vu_1$, iterating with the
  projected operator $\mat{P}\mat{A}\mat{P}$ (where
  $\mat{P}=\mat{I}-\vu_1\vu_1^{\!\top}$) produces the same sequence of directions as
  iterating with the deflated operator $\mat{A}-\mu_1\vu_1\vu_1^{\!\top}$, when both
  start from a vector already orthogonal to $\vu_1$. Why does the projection form
  not require knowing $\mu_1$?
  \item \textbf{Why orthogonality is mandatory.} Suppose you deflate with a stored
  eigenvector $\tilde\vu_1$ that is slightly \emph{inaccurate}---off the true
  $\vu_1$ by a small angle. Argue that the deflated iteration will then slowly
  re-grow a component of $\vu_1$ and can drift back toward the first eigenpair.
  What does this say about how accurately a pair must converge before it is safe to
  lock? (This is the practical reason locking tolerances exist in Krylov--Schur,
  \cref{ch:krylovschur}.)
  \item \textbf{Scalar Newton-on-$\lambda$.} Repeat \cref{wex:nepscalar} for
  $T(\lambda)=\lambda^2-(3+\sin\lambda)$ starting from $\lambda_0=2$. Carry three
  Newton steps and tabulate $\lambda_k$ and $\abs{T(\lambda_k)}$. Confirm the
  residual roughly squares each step (quadratic convergence). What would
  \emph{linear} convergence at rate $\tfrac14$ look like in the same table?
  \item \textbf{The normalization equation.} The NEP system
  \cref{eq:nep-system} appends $\vc^{\!\top}\vx=1$ to make the problem square.
  Explain why \emph{some} normalization is necessary (what goes wrong with the
  $N$ equations $\mat{T}(\lambda)\vx=\mathbf 0$ alone?). Does the choice of $\vc$
  affect \emph{which} eigenpair Newton converges to, or only the \emph{scale} of
  the eigenvector it returns? Give a sentence of reasoning for each.
  \item \textbf{Linearization's limit.} The quadratic pencil
  $(\Kop+\lambda\Cop+\lambda^2\Mop)\vx=\mathbf 0$ linearizes to a $2N\times2N$
  problem. A cubic pencil (degree three in $\lambda$) linearizes to what size? Now
  consider $\mat{T}(\lambda)=\Kop+\Mop\,e^{-\lambda\tau}$ (a delay term, genuinely
  transcendental). Explain why \emph{no} finite linearization exists, and what
  this forces the solver to do instead.
  \item \textbf{Dense-small versus sparse-large.} In the bordered solve of
  \cref{fig:dense-sparse-split}, the large block is $n\times n$ with
  $n\sim10^6$ and the small block is $k\times k$ with $k\sim50$. Estimate the
  ratio of arithmetic cost between a single dense LU factorization of the small
  block ($\sim\tfrac23 k^3$ operations) and a single iterative solve of the large
  block (take it as $\sim100$ matrix-vector products at $\sim10n$ operations each).
  Which dominates, and by how many orders of magnitude? What does your answer say
  about whether deflation ``slows down'' each Newton step?
  \item \textbf{Block deflation.} The chapter mostly deflates one pair at a time,
  but \cref{sec:def-locking} noted a block form that locks several at once into
  $\mat{X}$. Sketch how the extended problem \cref{eq:nep-extended} and its growth
  rule (append a column to $\mat{X}$, border $\mat{H}$) would change if the inner
  solver converged a \emph{block} of $b$ pairs simultaneously instead of one. What
  property of the nonlinear deflation makes one-at-a-time the natural default for
  the NEP, even though linear Krylov--Schur deflates in blocks? (Hint: re-read the
  note on minimality index one and the dependence of each pair's deflation on the
  prior pair being fully converged.)
\end{enumerate}
